\ifdefined\gkver\else\def\gkver{student}\fi
\documentclass[\gkver]{gaokaozhenti}
\newcommand{\ccnum}[1]{{\fontspec{Noto Sans CJK JP}#1}}
\begin{document}

\chapter{2007年全国理综Ⅱ}
%% paperType: 理综物理部分

\begin{enumerate}
\item
%% number: 14
%% paperName: 2007年全国理综Ⅱ第 14 题 6 分
%% typeId: 1
%% type: 单选题
%% chapter: 气体性质
%% point: 热力学第一定律
%% method: 无
%% score: 6
%% degree: 650
%% duplicateId: 0
%% body:
对一定质量的气体，下列说法正确的是\xzanswer{A}

\fourchoices[answer=A]
{在体积缓慢不断增大的过程中，气体一定对外界做功}
{在压强不断增大有过程中，外界对气体一定做功}
{在体积不断被压缩的过程中，内能一定增加}
{在与外界没有发生热量交换的过程中，内能一定不变}

%% answer: A
%% memo:
\memoanswer{
当气体增大时，气体对外界做功，当气体减小时，外界对气体做功，故 A 选项正确；根据 $\frac{pV}{T}=$ 常数，$p$ 增大时，$V$ 不一定变化，故 B 选项错误；在 $V$ 减小的过程中，可能向外界放热，根据 $\Delta E=W+Q$ 可知，内能不一定增大，故 C 选项错误；$Q=0$ 的过程中，$W$ 不一定为 0，故 D 选项错误。
}

\item
%% number: 15
%% paperName: 2007年全国理综Ⅱ第 15 题 6 分
%% typeId: 2
%% type: 多选题
%% chapter: 机械振动
%% point: 振动图象
%% method: 无
%% score: 6
%% degree: 650
%% duplicateId: 0
%% body:
一列横波在 $x$ 轴上传播，在 $x=0$ 与 $x=1\Ucm$ 的两点的振动图线分别如图中实线与虚线所示。由此可以得出\xzanswer{BC}

\onepicture[w=7.0cm]{figs/15.png}

\fourchoices[answer=BC]
{波长一定是 $4\Ucm$}
{波的周期一定是 $4\Us$}
{波的振幅一定是 $2\Ucm$}
{波的传播速度一定是 $1\Ucms$}

%% answer: BC
%% memo:
\memoanswer{
根据振动图象两个最大值的横坐标之差为振动周期，故 $T=4\Us$，B 选项正确；从图象可看出振幅 $A=2\Ucm$，C 选项正确；根据题中所给的振动图象无法得到波长（或波速），也就无法根据 $\lambda=v/T$ 算出波速（或波长），故 A、D 选项错误。
}

\item
%% number: 16
%% paperName: 2007年全国理综Ⅱ第 16 题 6 分
%% typeId: 1
%% type: 单选题
%% chapter: 运动和力
%% point: 动量和冲量
%% method: 无
%% score: 6
%% degree: 450
%% duplicateId: 0
%% body:
如图所示，PQS 是固定于竖直平面内的光滑的 $1/4$ 圆周轨道，圆心 O 在 S 的正上方。在 O、P 两点各有一质量为 $m$ 的有物块 a 和 b，从同时刻开始，a 自由下落，b 沿圆弧下滑。以下说法正确的是\xzanswer{A}

\onepicture[w=3.0cm]{figs/16.png}

\fourchoices[answer=A]
{a 比 b 先到达 S，它们在 S 点的动量不相等}
{a 与 b 先到达 S，它们在 S 点的动量不相等}
{a 比 b 先到达 S，它们在 S 点的动量相等}
{b 比 a 先到达 S，它们在 S 点的动量相等}

%% answer: A
%% memo:
\memoanswer{
a 做自由落体运动，$t_{a}=\sqrt{\frac{2R}{g}}$，b 做变速圆周运动，视为振动处理，$t_{b}=T/4=\sqrt{\frac{\pi^{2}R}{4g}}$，$t_{a}<t_{b}$；根据机械能守恒得知，a、b 到达 S 的速度大小相等，但方向不同，动量 $mv$ 大小相等，但方向不同，故 A 选项正确。
}

\item
%% number: 17
%% paperName: 2007年全国理综Ⅱ第 17 题 6 分
%% typeId: 2
%% type: 多选题
%% chapter: 光学
%% point: 光的折射
%% method: 无
%% score: 6
%% degree: 650
%% duplicateId: 0
%% body:
如图，P 是偏振片，P 的透振方向（用带的箭头的实线表示）为竖直方向。下列四种入射光束中，哪几种照射 P 时能在 P 的另一侧观察到透射光？\xzanswer{ABD}

\onepicture[w=6.0cm]{figs/17.png}

\fourchoices[answer=ABD]
{太阳光}
{沿竖直方向振动的光}
{沿水平方向振动的光}
{沿与竖直方向成 $45^{\circ}$ 角振动的光}

%% answer: ABD
%% memo:
\memoanswer{
根据光的现象，只要光的振动方向不与偏振片的狭缝垂直，光都能通过偏振片，太阳光沿后垂直传播方向的各个振动，能，沿竖直方向振动的光能通过偏振片，沿与竖直方向成 $45^{\circ}$ 角振动的光也能通过偏振片。故 A、B、D 都正确。
}

\item
%% number: 18
%% paperName: 2007年全国理综Ⅱ第 18 题 6 分
%% typeId: 2
%% type: 多选题
%% chapter: 原子和原子核
%% point: 玻尔模型
%% method: 无
%% score: 6
%% degree: 650
%% duplicateId: 0
%% body:
氢原子在某三个相邻能级间跃迁时，可发出三种不同波长的辐射光。已知其中的两个波长分别为 $\lambda_{1}$ 和 $\lambda_{2}$，且 $\lambda_{1}>\lambda_{2}$，则另一个波长可能是\xzanswer{CD}

\fourchoices[answer=CD]
{$\lambda_{1}+\lambda_{2}$}
{$\lambda_{1}-\lambda_{2}$}
{$\frac{\lambda_{1}\lambda_{2}}{\lambda_{1}+\lambda_{2}}$}
{$\frac{\lambda_{1}\lambda_{2}}{\lambda_{1}-\lambda_{2}}$}

%% answer: CD
%% memo:
\memoanswer{
玻尔原子模型的跃迁假设（$E_{\text{初}}-E_{\text{终}}=h\nu$）及 $\lambda=\frac{c}{\nu}$ 可得：$E_{3}-E_{1}=\frac{hc}{\lambda_{3}}$，$E_{3}-E_{2}=\frac{hc}{\lambda_{2}}$，$E_{2}-E_{1}=\frac{hc}{\lambda_{1}}$，所以得：$\frac{hc}{\lambda_{3}}=\frac{hc}{\lambda_{1}}+\frac{hc}{\lambda_{2}}\Rightarrow\lambda_{3}=\frac{\lambda_{1}\lambda_{2}}{\lambda_{1}+\lambda_{2}}$。

故 C 选项正确，同理 D 选项正确。

\onepicture[w=4.5cm]{figs/18_an.png}
}

\item
%% number: 19
%% paperName: 2007年全国理综Ⅱ第 19 题 6 分
%% typeId: 2
%% type: 多选题
%% chapter: 磁场
%% point: 带电粒子在磁场中的运动
%% method: 无
%% score: 6
%% degree: 650
%% duplicateId: 0
%% body:
如图所示，一带负电的质点在固定的正的点电荷作用下绕该正电荷做匀速圆周运动，周期为 $T_{0}$，轨道平面位于纸面内，质点的速度方向如图中箭头所示。现加一垂直于轨道平面的匀强磁场，已知轨道半径并不因此而改变，则\xzanswer{AD}

\onepicture[w=3.4cm]{figs/19.png}

\fourchoices[answer=AD]
{若磁场方向指向纸里，质点运动的周期大于 $T_{0}$}
{若磁场方向指向纸里，质点运动的周期小于 $T_{0}$}
{若磁场方向指向纸外，质点运动的周期小于 $T_{0}$}
{若磁场方向指向纸外，质点运动的周期小于 $T_{0}$}

%% answer: AD
%% memo:
\memoanswer{
匀速圆周运动、库仑定律、洛伦兹力、左手定则等知识列出：

未加磁场：$\frac{kQq}{r^{2}}=m\frac{4\pi^{2}}{T_{0}^{2}}r$　（1）

磁场指向纸里：$\frac{kQq}{r^{2}}-qvB=m\frac{4\pi^{2}}{T_{1}^{2}}r$　（2）

磁场指向纸外：$\frac{kQq}{r^{2}}+qvB=m\frac{4\pi^{2}}{T_{2}^{2}}r$　（3）

比较上述式子，$T_{1}>T_{0}$，$T_{2}<T_{0}$，故 AD 选项正确。
}

\item
%% number: 20
%% paperName: 2007年全国理综Ⅱ第 20 题 6 分
%% typeId: 2
%% type: 多选题
%% chapter: 功和能
%% point: 动能定理
%% method: 无
%% score: 6
%% degree: 450
%% duplicateId: 0
%% body:
假定地球、月球都静止不动，用火箭从地球沿地月连线发射一探测器。假定探测器在地球表面附近脱离火箭。用 $W$ 表示探测器从脱离火箭处到月球的过程中克服地球引力做的功，用 $E_{k}$ 表示探测器脱离火箭时的动能，若不计空气阻力，则\xzanswer{BD}

\fourchoices[answer=BD]
{$E_{k}$ 必须大于或等于 $W$，探测器才能到达月球}
{$E_{k}$ 小于 $W$，探测器也可能到达月球}
{$E_{k}=W/2$，探测器一定能到达月球}
{$E_{k}=W/2$，探测器一定不能到达月球}

%% answer: BD
%% memo:
\memoanswer{
设月球引力对探测器做的功为 $W_{1}$，根据动能定理可得：$-W+W_{1}=0-E_{k}$，根据 $F=G\frac{m_{1}m_{2}}{r^{2}}$ 可知，$F_{\text{地}}>F_{\text{月}}$，$W>W_{1}$，故 BD 选项正确。
}

\item
%% number: 21
%% paperName: 2007年全国理综Ⅱ第 21 题 6 分
%% typeId: 1
%% type: 单选题
%% chapter: 电磁感应
%% point: 电磁感应中图象
%% method: 无
%% score: 6
%% degree: 450
%% duplicateId: 0
%% body:
如图所示，在 PQ、QR 区域是在在着磁感应强度大小相等、方向相反的匀强磁场，磁场方向均垂直于纸面，bc 边与磁场的边界 P 重合。导线框与磁场区域的尺寸如图所示。从 $t=0$ 时刻开始线框匀速横穿两个磁场区域。以 a$\to$b$\to$c$\to$d$\to$e$\to$f 为线框中有电动势的正方向。以下四个 $\varepsilon-t$ 关系示意图中正确的是\xzanswer{C}

\onepicture[w=4.6cm]{figs/21.png}

\fourchoices[answer=C, ispicture=true, h=1.8cm]
{figs/21a.png}
{figs/21b.png}
{figs/21c.png}
{figs/21d.png}

%% answer: C
%% memo:
\memoanswer{
楞次定律或左手定则可判定线框刚开始进入磁场时，电流方向，即感应电动势的方向为顺时针方向，故 D 选项错误；$1\sim 2\Us$ 内，磁通量不变化，感应电动势为 0，A 选项错误；$2\sim 3\Us$ 内，产生感应电动势 $E=2Blv+Blv=3Blv$，感应电动势的方向为逆时针方向（正方向），故 C 选项正确。
}

\item
%% number: 22
%% paperName: 2007年全国理综Ⅱ第 22 题 5 分
%% typeId: 4
%% type: 实验题
%% chapter: 机械振动
%% point: 单摆
%% method: 无
%% score: 5
%% degree: 650
%% duplicateId: 0
%% body:
\begin{enumerate}
	\item
	在做“用单摆测定重力加速度”的实验中，有人提出以下几点建议：

	（A）适当加长摆线

	（B）质量相同，体积不同的摆球，应选用体积较大的

	（C）单摆偏离平衡位置的角度不能太大

	（D）单摆偏离平衡位置时开始计时，经过一次全振动后停止计时，用此时间间隔作为单摆摆动的周期

	其中对提高测量结果精度有利的是 \tkanswer[1.5cm]{AC}。

	\item
	有一电流表 A，量程为 $1\UmA$，内阻 $r_{g}$ 约为 $100\UO$。要求测量其内阻。可选用的器材有：电阻箱 $R_{0}$，最大阻值为 $99999.9\UO$；滑动变阻器甲，最大阻值为 $10\UkO$；滑动变阻器乙，最大阻值为 $2\UkO$；电源 $E_{1}$，电动势约为 $2\UV$，内阻不计；电源 $E_{2}$，电动势约为 $6\UV$，内阻不计；开关 2 个，导线若干。采用的测量电路图如图所示，实验步骤如下：

	（A）断开 $S_{1}$ 和 $S_{2}$，将 $R$ 调到最大

	（B）合上 $S_{1}$，调节 $R$ 使 A 满偏

	（C）合上 $S_{2}$，调节 $R_{1}$ 使 A 半偏，此时可认为的 A 的内阻 $r_{g}=R_{1}$。试问：

	（1）在上述可供选择的器材中，可变电阻 $R_{1}$ 应该选 \tkanswer[2cm]{$R_{0}$}；为了使测量尽是精确，可变电阻 $R$ 应该选择 \tkanswer[2cm]{滑动变阻器甲}；电源 $E$ 应该选择 \tkanswer[2cm]{$E_{2}$}。

	（2）认为内阻 $r_{g}=R_{1}$，此结果与的 $r_{g}$ 真实值相比 \tkanswer[1.5cm]{偏小}。（填“偏大”、“偏小”、“相等”）
\end{enumerate}

\onepicture[w=6.5cm, align=right]{figs/22.png}

%% answer:
\jdanswer{
\begin{enumerate}
	\item
	AC

	\item
	（1）$R_{0}$，滑动变阻器甲，$E_{2}$；（2）偏小
\end{enumerate}
}

%% memo:
\memoanswer{
（1）根据单摆的周期公式 $T=2\pi\sqrt{\frac{l}{g}}$ 可得，$g=\frac{4\pi^{2}l}{T^{2}}$，从该公式可看出，增大摆长 $l$，有利于减小误差，提高测量结果精度；$T$ 对测量结果影响较大，采用累计法测量以减小误差，故 C 无法提高测量结果精度；对 B 来说，由于球体积较大，空气阻力也大，单摆振动次数少，不利于采用累计法测量周期；故 B 不利于提高测量结果精度；只有在小角度的情形下，单摆的周期才满足 $T=2\pi\sqrt{\frac{l}{g}}$。综合上述，应选择 AC。

（2）（i）根据半偏法的测量原理，$R_{1}$ 必须选 $R_{0}$；由于电流表的满偏电流很小，要求 $R_{1}$ 的阻值很大，故 $R$ 应选滑动变阻器甲，电源选择 $E_{2}$，误差较小。

（ii）根据闭合电路的欧姆定律及电路特点可得：

合上 $S_{1}$，调节 $R$ 使 A 满偏：$I_{g}=\frac{E}{R+r_{g}+r}$

合上 $S_{2}$，调节 $R_{1}$ 使 A 半偏（电路中的总电流）：$I=\frac{E}{R+\frac{r_{g}R_{1}}{r_{g}+R_{1}}+r}$

故：$I>I_{g}$。

所以，通过电阻箱的电流 $I_{R1}$：$I_{R1}>I_{g}$

用 $U$ 表示电阻箱两端电压，则 $R_{1}=\frac{U}{I_{R1}}<\frac{U}{\frac{I_{g}}{2}}=I_{g}$

即：$r_{g}>R_{1}$

故：认为内阻 $r_{g}=R_{1}$，此结果与的 $r_{g}$ 真实值相比偏小。
}

\item
%% number: 23
%% paperName: 2007年全国理综Ⅱ第 23 题 16 分
%% typeId: 6
%% type: 计算题
%% chapter: 曲线运动
%% point: 竖直平面圆周运动
%% method: 无
%% score: 16
%% degree: 450
%% duplicateId: 0
%% body:
如图所示，位于竖直平面内的光滑轨道，由一段斜的直轨道的与之相切的圆形轨道连接而成，圆形轨道的半径为 $R$。一质量为 $m$ 的小物块从斜轨道上某处由静止开始下滑，然后沿圆形轨道运动。要求物块能通过圆形轨道的最高点，且在该最高点与轨道间的压力不能超过 $5mg$（$g$ 为重力加速度）。求物块初始位置相对于圆形轨道底部的高度 $h$ 的取值范围。

\onepicture[w=6.5cm, align=right]{figs/23.png}

%% answer:
\jdanswer{
$\frac{5}{2}R\leqslant h\leqslant 5R$
}

%% memo:
\memoanswer{
设物块在圆形轨道最高点的速度为 $v$，由机械能守恒得

$mgh=2mgR+\frac{1}{2}mv^{2}$　①

物块在最高点受的力为重力 $mg$、轨道的压力 $N$。重力与压力的合力提供向心力，有

$mg+N=m\frac{v^{2}}{R}$　②

物块能通过最高点的条件是

$N>0$　③

由②③式得

$v\geqslant \sqrt{gR}$　④

由①④式得

$h\geqslant \frac{5}{2}R$　⑤

按题的要求，$N\leqslant 5mg$，由②⑤式得

$v\leqslant \sqrt{6gR}$　⑥

由①⑥式得

$h\leqslant 5R$　⑦

$h$ 的取值范围是

$\frac{5}{2}R\leqslant h\leqslant 5R$
}

\item
%% number: 24
%% paperName: 2007年全国理综Ⅱ第 24 题 19 分
%% typeId: 6
%% type: 计算题
%% chapter: 运动和力
%% point: 动量和能量
%% method: 无
%% score: 19
%% degree: 450
%% duplicateId: 0
%% body:
用放射源钋的 $\alpha$ 射线轰击铍时，能发射出一种穿透力极强的中性射线，这就是所谓铍“辐射”。1932年，查德威克用铍“辐射”分别照射（轰击）氢和氮（它们可视为处于静止状态），测得照射后沿铍：“辐射”方向高速运动的氢核和氮核的速度之比为 7.0。查德威克假设铍“辐射”是由一种质量不为零的中性粒子构成的，从而通过实验在历史上首次发现了中子。假定铍“辐射”中的中性粒子与氢核或氮核发生弹性正碰，试在一考虑相对论效应的条件下计算构成铍“辐射”的中性粒子的质量。（质量用原子质量单位 $\Uu$ 表示，$1\Uu$ 等于 \ce{^{12}C} 原子质量的十二分之一。取氢核和氮核的质量分别为 $1.0\Uu$ 和 $14.0\Uu$。）

%% answer:
\jdanswer{
$m=1.2\Uu$
}

%% memo:
\memoanswer{
设构成铍“辐射”的中性粒子的质量和速度分别为 $m$ 和 $v$，氢核的质量为 $m_{H}$。构成铍“辐射”的中性粒子与氢核发生弹性正碰，碰后两粒子的速度分别为 $v'$ 和 $v_{H}'$。由动量守恒与能量守恒定律得

$mv=mv'+m_{H}v_{H}'$　①

$\frac{1}{2}mv^{2}=\frac{1}{2}mv'^{2}+\frac{1}{2}m_{H}v_{H}'^{2}$　②

解得

$v_{H}'=\frac{2mv}{m+m_{H}}$　③

同理，对于质量为 $m_{N}$ 的氮核，其碰后速度为

$v_{N}'=\frac{2mv}{m+m_{N}}$　④

由③④式可得

$m=\frac{m_{N}v_{N}'-m_{H}v_{H}'}{v_{H}'-v_{N}'}$　⑤

根据题意可知

$v_{H}'=7.0v_{N}'$　⑥

将上式与题给数据代入⑤式得

$m=1.2\Uu$　⑦
}

\item
%% number: 25
%% paperName: 2007年全国理综Ⅱ第 25 题 20 分
%% typeId: 6
%% type: 计算题
%% chapter: 磁场
%% point: 带电粒子在磁场中的运动
%% method: 无
%% score: 20
%% degree: 250
%% duplicateId: 0
%% body:
如图所示，在坐标系 $Oxy$ 的第一象限中在在沿 $y$ 轴正方向的匀强电场，场强大小为 $E$。在其它象限中在在匀强磁场，磁场方向垂直于纸面向里，A 是 $y$ 轴上的一点，它到坐标原点 O 的距离为 $h$；C 是 $x$ 轴上的一点，到 O 点的距离为 $l$，一质量为 $m$、电荷量为 $q$ 的带负电的粒子以某一初速度沿 $x$ 轴方向从 A 点进入电场区域，继而通过 C 点进入磁场区域，并再次通过 A 点，此时速度方向与 $y$ 轴正方向成锐角。不计重力作用。试求：

\begin{enumerate}
	\item
	粒子经过 C 点时速度的大小和方向；
	\item
	磁感应强度的大小 $B$。
\end{enumerate}

\onepicture[w=5.0cm, align=right]{figs/25.png}

%% answer:
\jdanswer{
\begin{enumerate}
	\item
	$v=\sqrt{\frac{qE(4h^{2}+l^{2})}{2mh}}$，方向与 $x$ 轴正方向的夹角 $\alpha=\arctan\frac{2h}{l}$

	\item
	$B=\frac{l}{h^{2}+l^{2}}\sqrt{\frac{2mhE}{q}}$
\end{enumerate}
}

%% memo:
\memoanswer{
（1）以 $a$ 表示粒子在电场作用下的加速度，有

$qE=ma$　①

加速度沿 $y$ 轴负方向。设粒子从 A 点进入电场时的初速度为 $v_{0}$，由 A 点运动到 C 点经历的时间为 $t$，则有

$h=\frac{1}{2}at^{2}$　②

$l=v_{0}t$　③

由②③式得

$v_{0}=l\sqrt{\frac{a}{2h}}$　④

设粒子从 C 点进入磁场时的速度为 $v$，$v$ 垂直于 $x$ 轴的分量

$v_{1}=\sqrt{2ah}$　⑤

由①④⑤式得

$v=\sqrt{v_{0}^{2}+v_{1}^{2}}=\sqrt{\frac{qE(4h^{2}+l^{2})}{2mh}}$　⑥

\onepicture[w=5.5cm]{figs/25_an.jpg}

设粒子经过 C 点时的速度方向与 $x$ 轴的夹角为 $\alpha$，则有

$\tan\alpha=\frac{v_{1}}{v_{0}}$　⑦

由④⑤⑦式得

$\alpha=\arctan\frac{2h}{l}$　⑧

（2）粒子从 C 点进入磁场后在磁场中作速度为 $v$ 的圆周运动。若圆周的半径为 $R$，则有

$qvB=m\frac{v}{R}$　⑨

设圆心为 P，则 PC 必与过 C 的速度垂直，且有 $\overrightarrow{PC}=\overrightarrow{PA}=R$。用 $\beta$ 表示 $\overrightarrow{PA}$ 与 $y$ 轴的夹角，由几何关系得

$R\cos\beta=R\cos\alpha+h$　⑩

$R\sin\beta=l-R\sin\alpha$　\ccnum{⑪}

由⑧⑩式解得

$R=\frac{h^{2}+l^{2}}{2hl}\sqrt{4h^{2}+l^{2}}$　\ccnum{⑫}

由⑥⑨式得

$B=\frac{l}{h^{2}+l^{2}}\sqrt{\frac{2mhE}{q}}$
}

\end{enumerate}

\end{document}
